\chapter{Functions, Arguments \& Scope}

A \textbf{Function} is a reusable block of code designed to perform a specific task. Functions break down large complex programs into modular, readable, and maintainable units (the DRY principle: Don't Repeat Yourself).

\section{Defining and Calling Functions}

Functions are defined using the `def` keyword, followed by the function name, parentheses `()`, and a colon `:`.

\begin{definitionbox}{Function Definition Syntax}
\begin{lstlisting}
def greet_user(name):
    """Docstring: Greets a user by name."""
    message = f"Hello, {name}! Welcome to Python."
    return message

# Function Call
result = greet_user("Alice")
print(result)
\end{lstlisting}
\end{definitionbox}

\begin{videocard}{IITM BS Lecture: Introduction to Functions}{https://www.youtube.com/watch?v=zDZRfWWetg0}
Official lecture defining `def`, parameters, `return` values, and docstrings.
\end{videocard}

\section{Types of Function Arguments}

Python functions support several types of parameters:
1. \textbf{Positional Arguments:} Passed in the exact order declared.
2. \textbf{Keyword Arguments:} Passed by explicitly naming parameter names (`name="Alice"`).
3. \textbf{Default Arguments:} Parameters with fallback default values if omitted by the caller.
4. \textbf{Variable-length Arguments (`*args` and `**kwargs`):}
   - `*args`: Collects extra positional arguments into a tuple.
   - `**kwargs`: Collects extra keyword arguments into a dictionary.

\begin{theorembox}{Argument Types Syntax}
\begin{lstlisting}
def calculate_price(price, discount=0.1, tax=0.05):
    return price * (1 - discount) * (1 + tax)

print(calculate_price(100))               # Uses defaults: 100 * 0.9 * 1.05 = 94.5
print(calculate_price(100, tax=0.10))     # Keyword arg overrides tax only

def custom_sum(*args):
    return sum(args)                      # args is (1, 2, 3, 4)

print(custom_sum(1, 2, 3, 4))             # 10
\end{lstlisting}
\end{theorembox}

\textbf{Data Science Relevance:} Machine Learning functions in scikit-learn heavily use default and keyword arguments (e.g., `RandomForestClassifier(n_estimators=100, max_depth=5)`).

\begin{videocard}{IITM BS Lecture: Types of Function Arguments}{https://www.youtube.com/watch?v=NqjnCY2qQWU}
Official lecture explaining positional, keyword, default, `*args`, and `**kwargs`.
\end{videocard}

\section{Scope of Variables (LEGB Rule)}

\textbf{Scope} determines where a variable is accessible in your code. Python resolves variable names using the \textbf{LEGB Rule}:

\begin{formulabox}{The LEGB Rule}
\begin{enumerate}
    \item \textbf{L (Local):} Variables defined inside the current function body.
    \item \textbf{E (Enclosing):} Variables in enclosing/outer functions (nested functions).
    \item \textbf{G (Global):} Variables defined at the top-level of the script/module.
    \item \textbf{B (Built-in):} Built-in names in Python (`len`, `sum`, `range`, etc.).
\end{enumerate}
\end{formulabox}

\begin{examplebox}{Global vs Local Scope}
\begin{lstlisting}
counter = 0  # Global Variable

def increment():
    global counter  # Explicitly modify global variable
    counter += 1

increment()
print(counter)  # Output: 1
\end{lstlisting}
\end{examplebox}

\begin{videocard}{IITM BS Lecture: Scope of Variables}{https://www.youtube.com/watch?v=4q5rGHfR-ic}
Official lecture explaining Local vs Global scope, `global` keyword, and namespace resolution.
\end{videocard}

\newpage
\section{Practice Exercises}
\textit{Detailed step-by-step solutions are available in the Appendix.}

\begin{enumerate}
    \item \textbf{Basic Function:} Write a function `is_even(n)` that returns `True` if an integer $n$ is even, and `False` otherwise.

    \item \textbf{Default Arguments:} Write a function `compute_power(base, exponent=2)` that returns $\text{base}^{\text{exponent}}$. Test it with `compute_power(5)` and `compute_power(2, 3)`.

    \item \textbf{Flexible Aggregation (`*args`):} Write a function `calculate_mean(*args)` that accepts any number of numeric arguments and returns their arithmetic mean. Handle the edge case of no arguments passed by returning `0.0`.

    \item \textbf{Variable Scope:} What will be the output of the following script? Explain why.
\begin{lstlisting}
x = 50

def modify_val():
    x = 100
    print("Inside:", x)

modify_val()
print("Outside:", x)
\end{lstlisting}

    \item \textbf{Data Science Application (Min-Max Scaling Function):} Write a modular function `min_max_scaler(data)` that takes a list of numbers, calculates its min and max, and returns a new list of min-max normalized values rounded to 3 decimal places.
\end{enumerate}